% ======================================================================
% Ontology of Continua -- Core 1.3
% Cross-K module: crossk_k10_k11.tex
% Cross-level relations between K10 and K11
% ======================================================================

\subsubsection{Межуровневая структура K10–K11}
\label{sec:crossk-k10-k11}

Переход $K_{10} \rightarrow K_{11}$ обозначает возникновение
\textbf{транс-метатеоретического континуума} из метатеоретического.
Если $K_{10}$ даёт функториальные отображения, металогику и категориальные
структуры, связывающие теоретические системы, то $K_{11}$ вводит:
\begin{itemize}
  \item универсальные классы операторов, действующие между несколькими
        мета-рамками,
  \item условия более высокого порядка для совместимости метамоделей,
  \item глобальные пороги транс-метасогласованности,
  \item потоки, перестраивающие целые семейства метамоделей,
  \item циклы, регулирующие саму эволюцию метауровня,
  \item слой ограничений, обеспечивающий, чтобы эволюция моделей по всем уровням
        $K_0\dots K_{10}$ оставалась согласованной в более общем контуре.
\end{itemize}

В Core~1.3 $K_{11}$ рассматривается как активный верхний уровень доктрины, а не
как нерешённый placeholder. Здесь суммируется его структурная роль в
поднятой межуровневой архитектуре.

% ----------------------------------------------------------------------
\subsubsection{1. Задействованные уровни и их роли}

\paragraph{K10 (метатеоретический континуум).}
$K_{10}$ предоставляет:
\begin{itemize}
  \item оси $A^{10}$ (функторы, мета-логика, категории совместимости),
  \item потенциалы $P^{10}$ (мета-согласованность, экспрессивность, универсальность),
  \item пороги $\Theta^{10}$ (категориальная согласованность,
        мета-логическая корректность),
  \item потоки $J^{10}$ (интерпретация, трансляция, реконструкция),
  \item циклы $C^{10}$ (мета-стабильность, циклы уточнения).
\end{itemize}

\paragraph{K11 (транс-метатеоретический континуум).}
$K_{11}$, в формулировке Core~1.3, вводит:
\begin{itemize}
  \item оси $A^{11}$, описывающие отношения между целыми классами
        мета-рамок,
  \item потенциалы $P^{11}$, измеряющие глобальную согласованность между
        мета-уровнями,
  \item пороги $\Theta^{11}$ для транс-мета согласованности и совместимости,
  \item потоки $J^{11}$, перестраивающие семейства метамоделей,
  \item циклы $C^{11}$, регулирующие эволюцию структур $K_{10}$,
  \item меру $k(K_{11})$ универсальной структурной согласованности.
\end{itemize}

$K_{11}$ служит уровнем ограничения для всех нижних уровней моделирования.

% ----------------------------------------------------------------------
\subsubsection{2. Наследуемые / совместимые оси и пороги}

\paragraph{Наследуемая структура.}
Мета-модели в $K_{10}$ становятся объектами; $K_{11}$ вводит морфизмы между
мета-морфизмами:
\[
\mathsf{MetaFunctors}(K_{10}) \rightarrow \mathsf{TransFunctors}(K_{11}).
\]

\paragraph{Новые оси.}
Оси в $A^{11}$ включают:
\begin{itemize}
  \item $A_{\mathrm{ultra}}$ — классы операторов, действующие на категории
        мета-моделей,
  \item $A_{\mathrm{compat}}^{(11)}$ — транс-мета отношения совместимости,
  \item $A_{\mathrm{reflect}}^{(11)}$ — оси более высокого порядка саморефлексии,
  \item $A_{\mathrm{limit}}$ — ограничения, задаваемые универсальной согласованностью.
\end{itemize}

\paragraph{Новые пороги.}
$\Theta^{11}$ включает:
\begin{itemize}
  \item \textbf{порог транс-мета-согласованности} — нарушение приводит к
        коллапсу $K_{11}$,
  \item \textbf{порог совместимости} — несовместимые мета-рамки не могут
        сосуществовать,
  \item \textbf{порог универсальной согласованности} — запрещает парадокс на
        транс-мета уровне,
  \item \textbf{порог экспрессивности} — минимальная представленческая
        достаточность по отношению к всем $K_0\dots K_{10}$.
\end{itemize}

Нарушение $\Theta^{10}$ предотвращает возникновение $K_{11}$.

% ----------------------------------------------------------------------
\subsubsection{3. Межуровневые потоки, циклы и напряжения}

\paragraph{Потоки.}
Потоки в $K_{11}$ расширяют мета-потоки:
\[
J^{11} = (J_{\mathrm{ultra}},\ J_{\mathrm{crossmeta}},\
         J_{\mathrm{lift2}},\ J_{\mathrm{constraint}}),
\]
где:
\begin{itemize}
  \item $J_{\mathrm{ultra}}$ — перестройка категорий метамоделей,
  \item $J_{\mathrm{crossmeta}}$ — трансляция между несовместимыми мета-уровнями,
  \item $J_{\mathrm{lift2}}$ — представление мета-рамки как часть $K_{11}$,
  \item $J_{\mathrm{constraint}}$ — наложение глобальных структурных ограничений.
\end{itemize}

\paragraph{Циклы.}
$C^{11}$ включает циклы более высокого порядка, такие как:
\[
C_{\mathrm{ultra}}:\
\text{metamodel} \rightarrow \text{lift} \rightarrow \text{constraint}
\rightarrow \text{refinement} \rightarrow \text{metamodel}.
\]

Эти циклы регулируют эволюцию мета-континуума.

\paragraph{Напряжение.}
Межуровневое напряжение:
\[
T_{10,11} =
f(\Theta^{11} - P^{11},\ A^{11} - A^{10},\
  J^{11} - J^{10},\ \Theta^{10} - \Theta^{11}).
\]

Избыточное напряжение коллапсирует $K_{11}$ к более простой метатеоретической
режиме.

% ----------------------------------------------------------------------
\subsubsection{4. Условия рождения и исчезновения между уровнями}

\paragraph{Рождение $K_{11}$.}
$K_{11}$ возникает, когда:
\[
T_{10} > \Theta_{10,\mathrm{dim}}
\quad\text{и}\quad
A^{11} \in M_{11} \setminus A(K_{10}),
\]
то есть когда различия между целыми классами мета-рамок становятся необходимыми.

Расширение:
\[
\Omega(K_{11}) = \Omega(K_{10}) \cup \Delta\Omega_{\mathrm{ultra}}.
\]

\paragraph{Распространение гибели.}
Если $\Theta^{10}$ нарушена, транс-мета-слой не может существовать.

Если пороги $\Theta^{11}$ нарушены (несовместимость функторов,
мета-парадокс), то $K_{11}$ коллапсирует, но $K_{10}$ сохраняется:
\[
\Omega(K_{11}) \rightarrow \varnothing.
\]

% ----------------------------------------------------------------------
\subsubsection{5. Концептуальные примеры}

\paragraph{Сопоставление теорий.}
$K_{11}$ может обеспечивать операцию более высокого порядка, которая
сопоставляет теории:
\[
\mathcal{F} : \mathsf{Mod}(T_1) \rightarrow \mathsf{Mod}(T_2).
\]

\paragraph{Интерпретация парадигм.}
Немоделируемые в $K_9$ несовместимые теории могут быть согласованы посредством
их подъёма в $K_{11}$.

\paragraph{Универсальные ограничения.}
$K_{11}$ обеспечивает глобальную согласованность по всем уровням моделирования:
\[
K_{11} : \{K_0,\dots,K_{10}\} \rightarrow \text{coherent envelope}.
\]

\paragraph{Случай коллапса.}
Если появляется транс-мета несогласованность:
\[
T_{10,11} > \Theta^{11}_{\mathrm{coh}},
\]
$K_{11}$ коллапсирует в $K_{10}$.

Следовательно, переход K10→K11 в сжатой форме можно описать как
возникновение универсальных операторов и ограничений, способных соотносить и
регулировать целые классы метатеоретических рамок.
